\section{Physical ansatz and equivalences}
\label{sec:ansatz}

\subsection{Scope and logical encoding}

The channel-first construction in the accepted target
specification~\cite{Targets} fixes a small physical encoding before
imposing tensor stationarity. The question addressed here is whether
the coefficients of that encoding can be compatible with blocking
one site into two. This is a necessary linear condition for the
renormalization fixed-point equations of
Cirac, P\'erez-Garc\'ia, Schuch, and Verstraete~\cite{CPSV}.
It precedes complete positivity, recovery, and canonical analysis.

The logical algebra, physical Hilbert space, and redundant state are
\begin{align}
	\mc A & =M_2(\C)\oplus\C,
	      &
	\mc H & =(\C^2\otimes\C^2)\oplus\C,
	      &
	\tau  & =\diag(1/3,2/3).
	\label{eq:spaces}
\end{align}
The logical trace is $\tr(X\oplus z)=\tr X+z$. The accepted encoding is
\begin{align}
	J_1(X\oplus z) & =(X\otimes\tau)\oplus z.
	\label{eq:encoding}
\end{align}
The last physical basis vector is denoted by $\ket4$.
If $P_0$ projects onto the four-dimensional physical summand, define
\begin{align}
	D_1(Y)
	 & =\tr_{\mathrm{red}}(P_0YP_0)\oplus\bra4Y\ket4,
	\nonumber                                         \\
	D_1J_1(X\oplus z)
	 & =X\tr\tau\oplus z=X\oplus z,
	\nonumber                                         \\
	\tr D_1(Y)
	 & =\tr(P_0Y)+\bra4Y\ket4=\tr Y.
	\label{eq:decoder}
\end{align}
Here the partial trace acts on the second qubit of the first summand.
Compression, partial trace, and preparation of $\tau$ are completely
positive (CP). Thus $J_1$ and $D_1$ are CP and trace preserving with
the indicated direct-sum trace.

The horizontal indices are pairs
$\alpha=(u,u')$ and $\beta=(v,v')$, with all component indices in
$\{0,1\}$. The horizontal MPDO bond dimension is therefore four.
For physical operator letters $M^{\alpha\beta}$, periodic contraction
and blocking mean
\begin{align}
	\rho_N(M)
	 & =\sum_{\alpha_1,\ldots,\alpha_N}
	M^{\alpha_1\alpha_2}\otimes\cdots\otimes
	M^{\alpha_N\alpha_1},
	\nonumber                                                \\
	G^{\alpha\beta}
	 & =\sum_\gamma M^{\alpha\gamma}\otimes M^{\gamma\beta}.
	\label{eq:contractions}
\end{align}
A physical fixed-point certificate requires completely positive,
trace-preserving (CPTP) maps on the complete physical operator spaces:
\begin{align}
	\mc T:\End(\mc H)            & \longrightarrow\End(\mc H^{\otimes2}),
	                             &
	\mc T(M^{\alpha\beta})       & =G^{\alpha\beta},
	\nonumber                                                             \\
	\mc S:\End(\mc H^{\otimes2}) & \longrightarrow\End(\mc H),
	                             &
	\mc S(G^{\alpha\beta})       & =M^{\alpha\beta}.
	\label{eq:rfp}
\end{align}
These are identities for every complex letter. Recovery of encoded
states alone does not establish them.

\subsection{Local purification and Gram ranks}

Let $A_0,A_1,C_0,C_1$ be $2\times2$ complex matrices. Define physical
vectors by
\begin{align}
	v_{uv}   & =(A_0)_{uv}\ket0+(A_1)_{uv}\ket1,
	\nonumber                                    \\
	a_{uv,1} & =\sqrt{1/3}\,v_{uv}\otimes\ket0,
	         &
	a_{uv,2} & =\sqrt{2/3}\,v_{uv}\otimes\ket1,
	\nonumber                                    \\
	a_{uv,3} & =(C_0)_{uv}\ket4,
	         &
	a_{uv,4} & =(C_1)_{uv}\ket4.
	\label{eq:purification}
\end{align}
The first two vectors lie in the four-dimensional physical summand.
The locally purified tensor is
\begin{align}
	M^{(u,u'),(v,v')}
	 & =\sum_{e=1}^4
	\ket{a_{uv,e}}\bra{a_{u'v',e}}
	=J_1(X_{(u,u'),(v,v')}),
	\nonumber                                         \\
	X_{(u,u'),(v,v')}
	 & =\sum_{j,k=0}^1
	    (A_j)_{uv}\overline{(A_k)_{u'v'}}\ket j\bra k
	\nonumber                                         \\
	 & \quad\oplus
	\left((C_0)_{uv}\overline{(C_0)_{u'v'}}
	                +(C_1)_{uv}\overline{(C_1)_{u'v'}}\right).
	\label{eq:letters}
\end{align}

To identify the Gram ranks, introduce
\begin{align}
	\ket a    & =\sum_{u,v}\ket{u,v}\otimes v_{uv},
	          &
	\ket{c_l} & =\sum_{u,v}(C_l)_{uv}\ket{u,v},
	\nonumber                                           \\
	B_0       & =\ket a\bra a,
	          &
	B_1       & =\ket{c_0}\bra{c_0}+\ket{c_1}\bra{c_1}.
	\label{eq:sector-grams}
\end{align}
After ordering physical summands, the full positive Gram operator is
\begin{align}
	B & =(B_0\otimes\tau)\oplus B_1,
	\nonumber                           \\
	\rank B
	  & =\rank B_0\,\rank\tau+\rank B_1
	=2\rank B_0+\rank B_1.
	\label{eq:gram-rank}
\end{align}
The stratum studied here has
\begin{align}
	\rank B_0=1,\qquad \rank B_1=2.
	\label{eq:gram-stratum}
\end{align}
Its minimal environment dimension is exactly four.

Conversely, every positive Gram operator compatible with
\cref{eq:encoding} and satisfying \cref{eq:gram-stratum} admits
\cref{eq:purification}. A positive rank-one $B_0$ has a vector factor
$\ket a$. A positive rank-two $B_1$ has two independent vector factors,
obtained from its two positive eigenvalues and corresponding
orthonormal eigenvectors. Resolving those factors in the basis
$\ket{u,v}$ gives $A_0,A_1,C_0,C_1$ as above. Thus the four-matrix
description does not omit any tensor in the specified Gram stratum.

Local purification proves positivity at every length, independently
of stationarity. For an environment string $\boldsymbol e$, put
\begin{align}
	\ket{\Psi_{\boldsymbol e}}
	 & =\sum_{u_1,\ldots,u_N}
	\bigotimes_{n=1}^N\ket{a_{u_nu_{n+1},e_n}},
	\qquad u_{N+1}=u_1.
\end{align}
Expanding the periodic contraction gives
\begin{align}
	\rho_N(M)
	 & =\sum_{\boldsymbol e,\boldsymbol u,\boldsymbol u'}
	\bigotimes_{n=1}^N
	\ket{a_{u_nu_{n+1},e_n}}
	\bra{a_{u'_nu'_{n+1},e_n}}
	\nonumber                                             \\
	 & =\sum_{\boldsymbol e}
	\ket{\Psi_{\boldsymbol e}}\bra{\Psi_{\boldsymbol e}}
	\geq0.
	\label{eq:positivity}
\end{align}
This positivity statement is not a fixed-point certificate.

\subsection{Allowed gauges}

A common virtual similarity acts as
\begin{align}
	K & \longmapsto SKS^{-1},
	\qquad
	K\in\{A_0,A_1,C_0,C_1\},
	\qquad S\in\operatorname{GL}_2(\C).
	\label{eq:virtual-gauge}
\end{align}
It induces similarity by $S\otimes\overline S$ on all horizontal
coordinate matrices. It preserves periodic contractions and the
existence of maps satisfying \cref{eq:rfp}. It also preserves the
Gram ranks, because its action on the virtual vector factors is
invertible.

A unitary change of the pair $(C_0,C_1)$ leaves the letters unchanged.
A unitary change of $(A_0,A_1)$ conjugates the logical $M_2$ sector,
and is implemented physically by $(U\otimes\Id_2)\oplus1$.
It preserves coefficient compatibility and channel feasibility.
These frame changes will be quotiented out.

In contrast, a unitary mixing all four matrices preserves the trace
transfer but generally changes the physical tensor. Such mixing
cannot be discarded as a gauge of the fixed encoding. The relative
position of the active pair and the scalar pair is part of the
search.
